\documentclass[10pt,a4paper]{article}
\usepackage[utf8]{inputenc}
\usepackage[spanish]{babel}
\usepackage{times}
\usepackage{hyperref} % Re-enable hyperref
\usepackage{enumitem}
\usepackage{tikz}       % For graphics (might still be useful for other things, or can remove if only for pgf-radar)
% \usepackage{pgf-radar}  % No longer using pgf-radar, using Matplotlib image
\usepackage{amsmath}
\usepackage{amsfonts}
\usepackage{amssymb}
\usepackage{graphicx}
\usepackage{geometry}
\geometry{a4paper, margin=1in}
\usepackage{array}
\usepackage{sectsty}
\usepackage{ragged2e}

% Hyperlink setup - re-enabled
\hypersetup{
   colorlinks=false,
   pdfborder={0 0 0},
   breaklinks=true % Allows links to break across lines
}

\sectionfont{\large\bfseries}
\subsectionfont{\normalsize\bfseries}

\begin{document}

% --- Title and Puntuacion ---
\begin{center}
    {\Huge\bfseries US beef: Hormone-treated food will not enter UK after US deal\par}
    \vspace{0.5em}
    {\large\bfseries Puntuación: 7\par}
\end{center}
\vspace{1em}

% --- Autor, Fuente, Fecha ---
\noindent\textbf{Autor:} N/A \hfill \textbf{Fuente:} Fuente no identificada \hfill \textbf{Fecha:} N/A\par
\vspace{1em}

% --- Cuerpo de la Noticia ---
\section*{Cuerpo de la Noticia}
\setlength{\parskip}{1em}
\setlength{\parindent}{0em}
Hormone-treated beef will not enter UK after US deal, says government
\par\medskip
The National Farmers' Union said it was asking the government to provide more details on how checks would work to ensure safety standards were maintained.
\par\medskip
However, the government said certification procedures and border checks would ensure hormone-reared beef would not enter the UK.
\par\medskip
Some farmers and consumers have expressed fears that the deal could open the door to beef from cattle raised using hormones to boost their growth.
\par\medskip
The government has insisted that American hormone-treated meat will not start to seep into the UK market, following the tariff deal agreed this week that boosts the trade in beef in both directions.
\par\medskip
"The rules on food standards have not changed and they will not change as a result of the deal," said Chief Secretary to the Treasury Darren Jones.
\par\medskip
He added that the agencies responsible for border safety checks would be able to test meat for traces of hormone with "consequences" for anyone breaking the law.
\par\medskip
But Ian McCubbine, a beef farmer in Surrey, said he was concerned about the prospect of more American-grown beef coming into the UK.
\par\medskip
"How do we know what they are putting in?" he said, speaking to the BBC's Today programme.
\par\medskip
"We spent 50 years building an industry that is strong on environmental gain and animal welfare.
\par\medskip
"The concern is that the US [beef imports] could be of lower quality."
\par\medskip
The UK stopped allowing hormone-produced beef in 1989, when the practice was banned across the EU after it declared it unsafe.
\par\medskip
But many American farmers use growth hormones as a standard part of their beef production. Adding growth hormones makes cows put on muscle mass, and so makes their beef cheaper.
\par\medskip
The US and other countries that use the method, including Australia, say there is no added health-risk from hormone-treated beef.
\par\medskip
However, a lot of consumers are wary of it, with some commenting online that they would look out for UK-produced beef in future.
\par\medskip
As part of the tariff deal the UK has agreed to allow up to 13,000 metric tonnes of beef imports from the US tariff-free. That is the equivalent of one medium-sized steak per adult per year. Currently the US exports around 1,000 tonnes to the UK with a 20\% tariff, the UK's Department for Environment and Rural Affairs (Defra) said.
\par\medskip
In exchange, the UK will also be able to sell more beef to the US than it currently does, also up to 13,000 tonnes at a lower tariff than at present.
\par\medskip
The deal also includes lower tariffs on UK-made cars destined for the US market and US ethanol exports to the UK.
\vspace{1em}

% --- URL (Re-enabled and moved) ---
\section*{URL}
\url{ https://www.bbc.com/news/articles/c89pw3j7z9zo }
\vspace{1em}

% --- Resumen de la Valoración ---
\section*{Resumen de la Valoración}
    La noticia es clara y equilibrada, con diversas perspectivas y declaraciones cre\'{i}bles, pero necesita mayor profundidad cient\'{i}fica y contextualizaci\'{o}n t\'{e}cnica.
\vspace{1em}

% --- Resumen de la Valoración del Titular ---
\section*{Resumen de la Valoración del Titular}
    Titular evaluado con alta calidad period\'{i}stica: veraz, claro, relevante y sin elementos sensacionalistas, cumpliendo est\'{a}ndares profesionales de comunicaci\'{o}n.
\vspace{1em}

% --- Valoracion General (String Summary) ---
\section*{Valoración General}
    La noticia "US beef: Hormone-treated food will not enter UK after US deal" ofrece una estructura clara y un equilibrio informativo al incluir perspectivas de diversas partes, como el gobierno, agricultores y consumidores. Se destacan las declaraciones de figuras relevantes, lo que a\~{n}ade credibilidad y autenticidad al relato. Sin embargo, se identifican oportunidades para enriquecer la cobertura mediante una mayor profundizaci\'{o}n en los riesgos cient\'{i}ficos asociados con la carne tratada con hormonas y una descripci\'{o}n m\'{a}s detallada de los mecanismos de verificaci\'{o}n fronteriza. Asimismo, la incorporaci\'{o}n de estudios cient\'{i}ficos recientes sobre la seguridad de la carne con hormonas y un an\'{a}lisis del contexto y las implicaciones del acuerdo comercial agregar\'{i}an valor. En suma, la noticia cubre los aspectos esenciales del acuerdo, pero una mayor especificidad y contextualizaci\'{o}n en aspectos t\'{e}cnicos y cient\'{i}ficos mejorar\'{i}an la comprensi\'{o}n completa del tema.
\vspace{1em}

% --- Valoración del Titular ---
\section*{Valoración del Titular}
        \textbf{Análisis del titular:} \\ 
        \textbf{Mejora de la Respuesta sobre el Titular:}
\par\medskip
1. \textbf{Concepto y funci\'{o}n del elemento}
\par\medskip
   El titular presenta de manera clara y concisa informaci\'{o}n fundamental sobre una restricci\'{o}n comercial: la carne de res tratada con hormonas de EE.UU. no ingresar\'{a} al Reino Unido debido a un nuevo acuerdo. Este captura la esencia de un control riguroso sobre la importaci\'{o}n de alimentos espec\'{i}ficos desde EE.UU., subrayando una inquietud por la seguridad alimentaria. Responde claramente a la pregunta "Qu\'{e} hace qu\'{e}", especificando una prohibici\'{o}n directa del comercio de carne de res tratada con hormonas. Se evidencia un beneficio impl\'{i}cito, ya que la medida parece destinada a proteger la salud de los consumidores brit\'{a}nicos. El impacto humano se centra en los beneficios potenciales para la salud y la tranquilidad de los consumidores en el Reino Unido.
\par\medskip
2. \textbf{Enfoque y contenido}
\par\medskip
   El titular comunica impl\'{i}citamente el beneficio de proteger a la salud p\'{u}blica al evitar la entrada de carne tratada con hormonas. Aunque no detalla el impacto en las personas, sugiere una ventaja clara para la salud de los ciudadanos. Proporciona un valor informativo significativo al abordar la tensi\'{o}n entre el comercio internacional y los est\'{a}ndares de salud, permitiendo una comprensi\'{o}n b\'{a}sica de las implicaciones para los consumidores brit\'{a}nicos.
\par\medskip
3. \textbf{Estructura}
\par\medskip
   La estructura del titular es correcta (sujeto: "US beef" + verbo: "will not enter" + predicado: "UK after US deal"). Presenta expl\'{i}citamente al sujeto principal (la carne de res de EE.UU.) y la acci\'{o}n principal (no entrar al Reino Unido). No omite elementos clave y mantiene un orden sint\'{a}ctico claro.
\par\medskip
4. \textbf{Estilo}
\par\medskip
   La brevedad del titular (nueve palabras) lo hace accesible y f\'{a}cil de entender para el p\'{u}blico general. Utiliza un lenguaje claro y preciso, sin tecnicismos innecesarios, y emplea adecuadamente las may\'{u}sculas para nombres propios (US, UK), reforzando su claridad.
\par\medskip
5. \textbf{Errores comunes detectados}
\par\medskip
   No se detectan errores tales como iniciar con un verbo, desorden sint\'{a}ctico, ausencia de sujeto o acci\'{o}n, ni formas de sensacionalismo o clickbait.
\par\medskip
\textbf{CONCLUSI\'{O}N GENERAL: Aprobada}
\par\medskip
El titular satisface adecuadamente los criterios period\'{i}sticos, ofreciendo una presentaci\'{o}n clara, concisa y relevante sobre una decisi\'{o}n importante en el \'{a}mbito de la seguridad alimentaria y el comercio internacional.
\par\medskip
\textbf{TITULO PROPUESTO:} Carne de res de EE.UU. tratada con hormonas excluida del Reino Unido tras acuerdo
        \vspace{1em}
        % Titular reformulado eliminado según petición del usuario
\vspace{1em}

% --- Texto Referencia (Consolidated) ---
\section*{Texto de Referencia}
    \textit{(Mostrando desde el campo de diccionario directo: texto\_referencia\_diccionario)}
    \begin{itemize}[leftmargin=*]
            \item \textbf{Referencia La noticia sobre el acuerdo del Reino Unido con Estados Unidos en relaci\'{o}n a la carne de res tratada con hormonas es adecuada en t\'{e}rminos generales, pues incluye declaraciones atribuibles a figuras reconocibles como Darren Jones, Ian McCubbine y la National Farmers' Union\_ Esto aporta un nivel importante de fiabilidad\_ La inclusi\'{o}n de cifras precisas, tales como la cuota de importaci\'{o}n de 13,000 toneladas m\'{e}tricas, tambi\'{e}n a\~{n}ade rigor a los reportes\_:} 1
            \item \textbf{Referencia Hormone-treated beef will not enter UK after US deal, says government\_ The National Farmers' Union said it was asking the government to provide more details on how checks would work to ensure safety standards were maintained\_ ``The rules on food standards have not changed and they will not change as a result of the deal,'' said Chief Secretary to the Treasury Darren Jones\_ But Ian McCubbine, a beef farmer in Surrey, said he was concerned about the prospect of more American-grown beef coming into the UK\_ As part of the tariff deal the UK has agreed to allow up to 13,000 metric tonnes of beef imports from the US tariff-free\_:} 1
            \item \textbf{Referencia La noticia titulada ``US beef: Hormone-treated food will not enter UK after US deal'' es adecuada y cumple con los est\'{a}ndares esperados al ofrecer datos necesarios y declaraciones relevantes\_ Se distingue claramente entre hechos, interpretaciones y opiniones, lo que facilita una comprensi\'{o}n clara por parte del lector\_ Esto se logra a trav\'{e}s de declaraciones directas de fuentes identificadas, como el gobierno brit\'{a}nico y el National Farmers' Union, lo cual asegura una distinci\'{o}n clara entre relato y hechos\_:} 2
            \item \textbf{Referencia Hormone-treated beef will not enter UK after US deal, says government\_ The National Farmers' Union said it was asking the government to provide more details on how checks would work to ensure safety standards were maintained\_ ``The rules on food standards have not changed and they will not change as a result of the deal,'' said Chief Secretary to the Treasury Darren Jones\_:} 2
            \item \textbf{Referencia La noticia titulada ``US beef: Hormone-treated food will not enter UK after US deal'' es clara y bien estructurada, pero ofrece oportunidades para una contextualizaci\'{o}n m\'{a}s rica\_ Aqu\'{i} presentamos un an\'{a}lisis detallado de sus puntos fuertes y \'{a}reas donde podr\'{i}a mejorar::} 3
            \item \textbf{Referencia Hormone-treated beef will not enter UK after US deal, says government\_ The government has insisted that American hormone-treated meat will not start to seep into the UK market, following the tariff deal agreed this week that boosts the trade in beef in both directions\_ ``The rules on food standards have not changed and they will not change as a result of the deal,'' said Chief Secretary to the Treasury Darren Jones\_:} 5
            \item \textbf{Referencia Para mejorar el an\'{a}lisis de la noticia ``US beef: Hormone-treated food will not enter UK after US deal'', se pueden considerar las siguientes \'{a}reas::} 4
            \item \textbf{Referencia Hormone-treated beef will not enter UK after US deal, says government\_ The UK stopped allowing hormone-produced beef in 1989, when the practice was banned across the EU after it declared it unsafe\_ But many American farmers use growth hormones as a standard part of their beef production\_ Adding growth hormones makes cows put on muscle mass, and so makes their beef cheaper\_ The US and other countries that use the method, including Australia, say there is no added health-risk from hormone-treated beef\_:} 4
            \item \textbf{Referencia La noticia sobre la garant\'{i}a de que la carne tratada con hormonas de EE\_UU\_ no ingresar\'{a} al Reino Unido proporciona datos y declaraciones relevantes, manteniendo un equilibrio informativo al incluir perspectivas del gobierno, agricultores y consumidores\_ Este enfoque es efectivo para resaltar el n\'{u}cleo del acontecimiento: la promesa gubernamental de bloquear la entrada de este tipo de carne al pa\'{i}s\_:} 5
    \end{itemize}
\vspace{1em}

% --- Valoraciones (Dictionary of AI interpretations) ---
% This section is to be removed as per user request.
% \section*{Valoraciones (Interpretaciones IA)}
% %    ... (content previously here) ...
% % \vspace{1em}

% --- Análisis de Verificación de Hechos ---
\section*{Análisis de Verificación de Hechos}
    \# Verificaci\'{o}n de Hechos: Noticia sobre Carne de Res Tratada con Hormonas en el Reino Unido
\par\medskip
\#\# An\'{a}lisis de Afirmaciones Clave
\par\medskip
La noticia contiene varias afirmaciones espec\'{i}ficas que requieren verificaci\'{o}n contra las fuentes disponibles. A continuaci\'{o}n se presenta un an\'{a}lisis detallado de los datos num\'{e}ricos y afirmaciones concretas presentadas en el texto.
\par\medskip
\#\#\# Cronolog\'{i}a de la Prohibici\'{o}n de Carne Tratada con Hormonas
\par\medskip
La noticia se\~{n}ala que "El Reino Unido dej\'{o} de permitir carne de res producida con hormonas en 1989, cuando la pr\'{a}ctica fue prohibida en toda la UE despu\'{e}s de que la declarara insegura". Esta afirmaci\'{o}n es precisa[1][2][7]. Sin embargo, es importante notar que una fuente indica que la prohibici\'{o}n para el ganado dom\'{e}stico en el Reino Unido en realidad se remonta a 1981, mientras que la importaci\'{o}n fue prohibida en 1989[1]. La noticia no menciona esta distinci\'{o}n importante sobre la prohibici\'{o}n anterior para la producci\'{o}n dom\'{e}stica, lo cual podr\'{i}a considerarse una omisi\'{o}n significativa pero no una inexactitud directa.
\par\medskip
\#\#\# Datos de Exportaci\'{o}n Actuales
\par\medskip
La noticia establece que "Actualmente, EE.UU. exporta alrededor de 1,000 toneladas al Reino Unido con un arancel del 20\%". Sin embargo, existe una discrepancia en los datos m\'{a}s recientes. Las fuentes indican que "el U.S. export\'{o} 1,970 toneladas de carne de res al Reino Unido valoradas en 32 millones de d\'{o}lares en 2024"[1]. La noticia utiliza la palabra "alrededor" lo que proporciona cierta flexibilidad, pero 1,970 toneladas es considerablemente diferente de 1,000 toneladas, representando casi el doble. Esto constituye una imprecisi\'{o}n significativa en los datos de volumen.
\par\medskip
Adicionalmente, se menciona que la cuota anterior era de 1,000 toneladas m\'{e}tricas dentro de una tarifa espec\'{i}fica, lo que podr\'{i}a explicar parcialmente la referencia, pero esto no se aclara en la noticia[4][28][49][51].
\par\medskip
\#\#\# T\'{e}rminos del Nuevo Acuerdo de Comercio
\par\medskip
La afirmaci\'{o}n de que "el Reino Unido ha acordado permitir hasta 13,000 toneladas m\'{e}tricas de importaciones de carne de res de EE.UU. sin aranceles" es correcta[1][2][4][22]. Las fuentes m\'{u}ltiples confirman que la cuota preferencial sin derechos es de 13,000 toneladas m\'{e}tricas[4][22][49][51].
\par\medskip
La noticia tambi\'{e}n se\~{n}ala correctamente que "En cambio, el Reino Unido tambi\'{e}n podr\'{a} vender m\'{a}s carne de res a EE.UU. de la que actualmente lo hace, tambi\'{e}n hasta 13,000 toneladas con una tarifa m\'{a}s baja"[1][2][4][22]. Las fuentes confirman esta disposici\'{o}n rec\'{i}proca.
\par\medskip
\#\#\# An\'{a}lisis de Equivalencia de Consumo
\par\medskip
La noticia afirma que 13,000 toneladas m\'{e}tricas equivale a "un bistec mediano por adulto al a\~{n}o". Aunque esta declaraci\'{o}n aparece en las fuentes[2][10], es importante evaluar si esta equivalencia es matem\'{a}ticamente precisa. Con una poblaci\'{o}n del Reino Unido de aproximadamente 67 millones de personas, 13,000 toneladas divididas por 67 millones resultar\'{i}a en aproximadamente 194 gramos por persona (asumiendo 1,000 gramos = 1 kilogramo y 1 tonelada = 1,000 kilogramos). Un bistec mediano t\'{i}picamente pesa entre 150-250 gramos, por lo que la equivalencia parece razonable, aunque es una generalizaci\'{o}n que simplifica el consumo.
\par\medskip
\#\#\# Aseveraciones sobre Mecanismos de Control
\par\medskip
La noticia cita al Secretario Jefe del Tesoro, Darren Jones, diciendo "Las reglas sobre normas alimentarias no han cambiado y no cambiar\'{a}n como resultado del acuerdo"[1][2][7][10][13][25][31]. Esta cita aparece con consistencia en m\'{u}ltiples fuentes, confirmando su precisi\'{o}n.
\par\medskip
Respecto a los procedimientos de certificaci\'{o}n y pruebas, la noticia se\~{n}ala que "los organismos responsables de las comprobaciones de seguridad fronteriza podr\'{i}an probar la carne para detectar rastros de hormonas con 'consecuencias' para cualquiera que incumpliera la ley"[1][2]. Sin embargo, es importante notar que una fuente experta indica que "Probar la presencia de hormonas es posible pero sumamente dif\'{i}cil y requiere equipos muy costosos, siendo que el costo por prueba ronda en cientos de libras"[16][20][38][41]. Esto sugiere que aunque el gobierno afirma tener capacidad de prueba, los expertos advierten sobre las dificultades pr\'{a}cticas y los costos significativos asociados con estas pruebas.
\par\medskip
\#\#\# Contexto del Desacuerdo Internacional
\par\medskip
La noticia menciona que "EE.UU. y otros pa\'{i}ses que utilizan este m\'{e}todo, incluida Australia, dicen que no hay riesgo sanitario adicional por la carne de res tratada con hormonas". Mientras que es cierto que EE.UU. defiende el uso de hormonas, este resumen simplifica un desacuerdo cient\'{i}fico m\'{a}s complejo. Las fuentes indican que existe un "desacuerdo significativo entre la UE y EE.UU. sobre las cuestiones de salud de la carne de res tratada con hormonas"[20][38][55]. La UE sostiene que uno de los hormones utilizados regularmente es cancer\'{i}geno, mientras que EE.UU. y Canad\'{a} argumentan ante la OMC que la evaluaci\'{o}n de riesgos de la UE es defectuosa[20][38][55].
\par\medskip
De manera m\'{a}s espec\'{i}fica, una evaluaci\'{o}n cient\'{i}fica de 1999 de la Comisi\'{o}n Europea concluy\'{o} que "el estradiol-17 tiene que considerarse un carcin\'{o}geno completo" que "ejerce efectos tanto iniciadores como promotores de tumores"[24][44]. Aunque esto refleja la posici\'{o}n de la UE, la noticia no captura esta complejidad cient\'{i}fica subyacente.
\par\medskip
\#\#\# Precisi\'{o}n Temporal del Acuerdo
\par\medskip
La noticia se refiere al "acuerdo comercial acordado esta semana" sin proporcionar una fecha espec\'{i}fica. Bas\'{a}ndome en los datos disponibles, el acuerdo de cooperaci\'{o}n econ\'{o}mica EE.UU.-Reino Unido fue anunciado el 8 de mayo de 2025[1][3][4][16][49][50][51], lo que sugiere que esta noticia fue publicada alrededor de esa fecha o poco despu\'{e}s (espec\'{i}ficamente el 9 de mayo de 2025 seg\'{u}n la informaci\'{o}n de fuente[2][10]).
\par\medskip
\#\#\# Afirmaciones sobre Preocupaciones de los Agricultores
\par\medskip
La cita de Ian McCubbine, ganadero de Sussex, es precisa seg\'{u}n se reporta en las fuentes[2][10][33][40]. Sin embargo, es notable que la Uni\'{o}n Nacional de Agricultores (NFU) requiri\'{o} aclaraciones sobre c\'{o}mo funcionar\'{i}an las comprobaciones[2][10][40]. Una fuente posterior del NFU, de enero de 2026, confirma que el gobierno ha "mantenido firme, salvaguardando nuestros sectores agr\'{i}colas m\'{a}s sensibles y manteniendo nuestros altos est\'{a}ndares de bienestar animal, ambientales y de seguridad alimentaria"[29][48], lo que sugiere que las garant\'{i}as proporcionadas fueron aceptadas, al menos formalmente.
\par\medskip
\#\#\# Aspectos Verificados pero Incompletamente Presentados
\par\medskip
La noticia no aborda completamente ciertos aspectos relevantes. Por ejemplo, no menciona que el acuerdo especifica que la carne de res importada debe ser "libre de hormonas"[4][19][22][28][36]. El acuerdo, de hecho, especifica una "cuota de 13,000t de carne de res libre de hormonas ringfenced"[4][22][36]. Esto es claramente compatible con la prohibici\'{o}n del Reino Unido, pero la noticia podr\'{i}a haber sido m\'{a}s expl\'{i}cita sobre este punto crucial.
\par\medskip
Tambi\'{e}n, la noticia no menciona que el documento del acuerdo "no constituye un acuerdo legalmente vinculante"[6][51], un detalle importante para entender el estado legal del arreglo.
\par\medskip
\#\# Conclusi\'{o}n de la Verificaci\'{o}n
\par\medskip
Los datos verificables en la noticia son generalmente correctos, con la excepci\'{o}n notable de la cifra de exportaciones actuales, donde la noticia indica "alrededor de 1,000 toneladas" cuando las fuentes m\'{a}s recientes disponibles muestran 1,970 toneladas en 2024. Las citas de funcionarios gubernamentales y actores de la industria se citan de manera precisa. Sin embargo, la noticia simplifica ciertos aspectos cient\'{i}ficos y t\'{e}cnicos complejos, particularmente respecto a las capacidades de prueba y el estado del desacuerdo cient\'{i}fico internacional sobre los riesgos de la carne de res tratada con hormonas. Estos no son errores factuales directos sino m\'{a}s bien omisiones que afectan la comprensi\'{o}n completa de la situaci\'{o}n.
\vspace{1em}

% --- Fuentes del Análisis de Verificación de Hechos ---
\section*{Fuentes del Análisis}
    \begin{enumerate}[leftmargin=*]
            \item \url{ https://www.dtnpf.com/agriculture/web/ag/blogs/ag-policy-blog/blog-post/2025/05/09/behind-us-uk-trade-deal-hormone-beef }
            \item \url{ https://www.cattlerange.com/articles/2025/05/hormone-treated-beef-will-not-enter-uk-after-us-trade-deal/ }
            \item \url{ https://www.venable.com/insights/publications/2025/05/america-firsts-first-trade-deal-white-house }
            \item \url{ https://www.nfuonline.com/news/us-reciprocal-beef-quota-in-place/ }
            \item \url{ https://www.soilassociation.org/causes-campaigns/top-10-risks-from-a-uk-us-trade-deal/what-is-hormone-treated-beef/ }
            \item \url{ https://www.cfr.org/articles/tracking-trumps-trade-deals }
            \item \url{ https://www.supplychainbrain.com/articles/41756-trade-deal-uk-will-not-allow-import-of-american-hormone-treated-beef }
            \item \url{ https://food.r-biopharm.com/analytes/residues/ }
            \item \url{ https://www.food.gov.uk/business-guidance/importing-products-of-animal-origin-poao }
            \item \url{ https://www.cattlerange.com/articles/2025/05/hormone-treated-beef-will-not-enter-uk-after-us-trade-deal/ }
            \item \url{ https://pmc.ncbi.nlm.nih.gov/articles/PMC3834504/ }
            \item \url{ https://www.gov.uk/guidance/beef-and-veal-marketing-standards }
            \item \url{ https://www.supplychainbrain.com/articles/41756-trade-deal-uk-will-not-allow-import-of-american-hormone-treated-beef }
            \item \url{ https://pmc.ncbi.nlm.nih.gov/articles/PMC9336673/ }
            \item \url{ https://www.gov.uk/guidance/importing-live-animals-or-animal-products-from-non-eu-countries }
            \item \url{ https://www.foodsafetynews.com/2025/05/uk-says-food-standards-protected-in-u-s-trade-deal/ }
            \item \url{ https://uknowledge.uky.edu/cgi/viewcontent.cgi?article=1013\&context=animalsci\_etds }
            \item \url{ https://www.food.gov.uk/business-guidance/chapter-3-imported-and-exported-meat-and-animals }
            \item \url{ https://www.farmersjournal.ie/beef/news/uk-gains-reciprocal-access-for-13-000t-of-beef-to-us-899333 }
            \item \url{ https://www.sciencemediacentre.org/expert-reaction-to-the-debate-on-hormone-treated-beef-and-chlorinated-chicken/ }
            \item \url{ https://www.fb.org/market-intel/wheres-the-hormone-free-beef }
            \item \url{ https://www.nfuonline.com/news/us-reciprocal-beef-quota-in-place/ }
            \item \url{ https://pmc.ncbi.nlm.nih.gov/articles/PMC3834504/ }
            \item \url{ https://pmc.ncbi.nlm.nih.gov/articles/PMC1115840/ }
            \item \url{ https://www.supplychainbrain.com/articles/41756-trade-deal-uk-will-not-allow-import-of-american-hormone-treated-beef }
            \item \url{ https://www.dtnpf.com/agriculture/web/ag/news/article/2025/05/09/behind-us-uk-trade-deal-hormone-beef }
            \item \url{ https://www.usa-beef.org/for-distributors/non-hormone-treated-cattle-programme/ }
            \item \url{ https://ielp.worldtradelaw.net/2025/05/the-general-terms-of-the-us-uk-trade-deal/ }
            \item \url{ https://www.farmersguide.co.uk/business/politics/government-pledges-food-standards-will-not-be-compromised-in-us-trade-talks/ }
            \item \url{ https://www.ams.usda.gov/services/imports-exports/nhtc }
            \item \url{ https://www.supplychainbrain.com/articles/41756-trade-deal-uk-will-not-allow-import-of-american-hormone-treated-beef }
            \item \url{ https://apps.fas.usda.gov/newgainapi/api/Report/DownloadReportByFileName?fileName=Livestock+and+Products+Annual\_The+Hague\_European+Union\_E42025-0004 }
            \item \url{ https://www.cattlerange.com/articles/2025/05/hormone-treated-beef-will-not-enter-uk-after-us-trade-deal/ }
            \item \url{ https://www.gov.uk/government/publications/uk-australia-fta-report-under-section-42-of-agriculture-act-2020/report-pursuant-to-section-42-of-the-agriculture-act-2020-web-version }
            \item \url{ https://www.gov.uk/government/statistics/united-kingdom-food-security-digest-2025/united-kingdom-food-security-digest-2025 }
            \item \url{ https://www.nfuonline.com/news/uk-us-trade-agreement/ }
            \item \url{ https://www.dtnpf.com/agriculture/web/ag/news/article/2025/05/09/behind-us-uk-trade-deal-hormone-beef }
            \item \url{ https://www.sciencemediacentre.org/expert-reaction-to-the-debate-on-hormone-treated-beef-and-chlorinated-chicken/ }
            \item \url{ https://www.gov.uk/government/publications/notice-to-traders-0126-usa-beef-preferential-tariff-quota }
            \item \url{ https://www.cattlerange.com/articles/2025/05/hormone-treated-beef-will-not-enter-uk-after-us-trade-deal/ }
            \item \url{ https://www.foodsafetynews.com/2025/05/uk-says-food-standards-protected-in-u-s-trade-deal/ }
            \item \url{ https://www.thebeefsite.com/news/eu-moves-to-tighten-import-checks-on-food-and-farm-goods }
            \item \url{ https://www.ams.usda.gov/services/imports-exports/nhtc }
            \item \url{ https://food.ec.europa.eu/document/download/f9e332dd-6cc7-4346-ba5a-3785969e256e\_en?filename=cs\_meat\_hormone-out50\_en.pdf\&prefLang=lt }
            \item \url{ https://www.supplychainbrain.com/articles/41756-trade-deal-uk-will-not-allow-import-of-american-hormone-treated-beef }
            \item \url{ https://www.usa-beef.org/for-distributors/non-hormone-treated-cattle-programme/ }
            \item \url{ https://pmc.ncbi.nlm.nih.gov/articles/PMC9336673/ }
            \item \url{ https://www.farmersguide.co.uk/business/politics/government-pledges-food-standards-will-not-be-compromised-in-us-trade-talks/ }
            \item \url{ https://ielp.worldtradelaw.net/2025/05/the-general-terms-of-the-us-uk-trade-deal/ }
            \item \url{ https://www.whitehouse.gov/fact-sheets/2025/05/fact-sheet-u-s-uk-reach-historic-trade-deal/ }
            \item \url{ https://www.gov.uk/government/publications/us-uk-economic-prosperity-deal-epd/general-terms-for-the-united-states-of-america-and-the-united-kingdom-of-great-britain-and-northern-ireland-economic-prosperity-deal-web-accessible-v }
            \item \url{ https://www.expanamarkets.com/insights/article/global-effects-of-new-us-and-uk-beef-deal/ }
            \item \url{ https://ustr.gov/about/policy-offices/press-office/fact-sheets/2025/may/fact-sheet-us-uk-reach-historic-trade-deal }
            \item \url{ https://www.whitehouse.gov/fact-sheets/2025/06/fact-sheet-implementing-the-general-terms-of-the-u-s-uk-economic-prosperity-deal/ }
            \item \url{ https://www.sciencemediacentre.org/expert-reaction-to-the-debate-on-hormone-treated-beef-and-chlorinated-chicken/ }
            \item \url{ https://www.thebeefsite.com/news/britain-starts-checks-on-fresh-food-imports-from-the-eu }
            \item \url{ https://blog.ansi.org/ansi/differences-between-eu-and-us-food-standards/ }
            \item \url{ https://phys.org/news/2020-11-uk-consumers-hormones-beef-chlorine.html }
            \item \url{ https://www.food.gov.uk/business-guidance/chapter-3-imported-and-exported-meat-and-animals }
            \item \url{ https://farrellymitchell.com/food-safety-assurance-standards/ }
    \end{enumerate}
\vspace{1em}

\end{document}